\documentclass[letterpaper]{article}

\usepackage[letterpaper,margin=1in]{geometry}

% Compile from the repository root so that this path resolves to
% the pinned WSP submodule.
\input{wsp/testing/test-case-library.tex}

\begin{document}

\def\TCID{%
  TC-0011
}

\def\TCTitle{%
  TabStrip and StatusBar
}

\def\TCPurpose{%
  Verify the TabStrip tabs, selection, layout metrics, hit testing and events,
  and the StatusBar panels, style and events.
}

\def\TCPriority{%
  \TCPriorityHigh
}

\def\TCRequirementRef{%
  REQ-0011, REQ-0017, REQ-0018, REQ-0019, REQ-0023, REQ-0024
}

\def\TCDesignRef{%
  \texttt{ADR-0006} and \texttt{ADR-0010} in WFC
}

\def\TCFeatureRef{%
  \texttt{tests/corpus/189-comctl-tabstatus.frm} and
  \texttt{tests/corpus/189-comctl-tabstatus.vbp}; the form code and the
  controls of \texttt{include/wfc/mscomctl/mscomctl.inc} and the runner
  \texttt{tests/tc-0011-comctl-tabstatus.ps1}
}

\def\TCTestTechnique{%
  Scenario-based testing of one form: each Immediate-window line is compared
  with the recorded outcome (ISO/IEC/IEEE 29119-4 experience-based and
  specification-based techniques, with expected output comparison).
}

\def\TCPreconditions{%
  A \texttt{wfc} of the same version as the library (WFC, the interpreter and
  compiler these controls belong to, runs the case: the library holds the
  program, WFC the runner). The files of the case are in
  \texttt{tests/corpus}: \texttt{189-comctl-tabstatus.vbp},
  \texttt{189-comctl-tabstatus.frm}, \texttt{189-comctl-tabstatus.expected},
  \texttt{189-comctl-tabstatus.debug}, \texttt{189-comctl-tabstatus.ui}.
}

\def\TCEnvironment{%
  Windows with \texttt{wfc.exe} of the x64 Debug build; the form runs in WFC's
  headless host (no desktop is needed). WFC registers the case with CTest and
  runs \texttt{wfc --interpret 189-comctl-tabstatus.vbp}.
}

\def\TCAssumptions{%
  The scripted input \texttt{189-comctl-tabstatus.ui} is the only input; the
  case does not read the registry or other files except the picture named in
  the form.
}

\def\TCInputData{%
  A form with a TabStrip and a StatusBar; tabs \texttt{zero}, \texttt{one},
  \texttt{two}; panels \texttt{msg} and \texttt{clock}; the script selects the
  tab \texttt{zero}, clicks panel 1 and double-clicks panel \texttt{clock}.
}

\def\TCInitialState{%
  A fresh \texttt{wfc} process; the program creates its own controls from the
  form file and starts with no state from a previous run.
}

\def\TCProcedure{%
  \begin{enumerate}
    \item Clear the tabs and panels, add three tabs (one at index 1) and print
      the count, keys, index, caption and type name.
    \item Print the selected tab, a tab's \texttt{Left}, \texttt{ClientTop}
      and \texttt{ClientLeft}, and \texttt{HitTest} on a tab and on an empty
      spot.
    \item Select a tab with \texttt{Set}, remove a tab, add a duplicate key,
      read an out-of-range index, and \texttt{DeselectAll}.
    \item Add two panels, print the panel count, text, style, \texttt{Left}
      and the bar style, and set and print \texttt{SimpleText}.
    \item Run the script and print the events.
  \end{enumerate}
}

\def\TCExpectedResult{%
  \begin{enumerate}
    \item Three tabs, the first key \texttt{zero}, \texttt{one} at index 2,
      caption \texttt{First}, type name \texttt{Tab}; the selected tab is
      \texttt{one}, the one added first, and \texttt{HitTest} returns the tab
      at its own point and \texttt{Nothing} at the bottom right corner.
    \item After the removal two tabs are left; a duplicate key is 35602, an
      index of 9 is 35600; after \texttt{DeselectAll} the selection is
      \texttt{Nothing}.
    \item Two panels, the first with text \texttt{Ready}, style 0 and a second
      panel to the right of it; \texttt{SimpleText} reads back
      \texttt{simple}.
    \item Events: \texttt{before} and \texttt{click Second} when the form
      selects a tab with \texttt{Set}; from the script \texttt{before},
      \texttt{click Zero}, \texttt{panel msg}, \texttt{panel clock} and
      \texttt{dbl clock}.
    \item \texttt{wfc} exits with code 0, writes the exact text of
      \texttt{189-comctl-tabstatus.expected} to standard output, and writes
      the exact Immediate-window text of \texttt{189-comctl-tabstatus.debug}
      to standard error. Any difference fails the case.
  \end{enumerate}
}

\def\TCPostConditions{%
  The form is closed by the last \texttt{close Form1} line of the script;
  nothing is left running and no file is written.
}

\TCTemplate

\end{document}
